% 第 41 章。数字与教学模型：checks/check_ch41.py；图：figsrc/fig_ch41.py。
% 史料：development_facts（取数脚本 data/fetch/development41.py）；六个经济体的人均 GDP：maddison2023_development（data/fetch/maddison2023_development41.py）。
\chapter{发展经济学：二元经济、进口替代与大推进}\label{ch:development}

\begin{chapterguide}
\guidequestion{穷国怎样起飞？}
\begin{guidegoals}
\item 推导刘易斯（1954）的二元经济模型，说明储蓄率为什么随资本主义部门的扩大而上升；
\item 概括普雷维什（1950）的贸易条件论，计算有效保护率，比较进口替代与出口导向；
\item 解释大推进的两个均衡与资金缺口模型，以及它们受到的检验；
\item 用印度与加纳的计划说明，独立后的发展政策由谁决定、代价由谁承担。
\end{guidegoals}
\guideprereq{\chref{ch:solow}\enspace\nameref{ch:solow}；\chref{ch:periphery}\enspace\nameref{ch:periphery}}
\end{chapterguide}

1952 年 8 月的一天，刘易斯（W. Arthur Lewis，1915—1991）走在曼谷的路上\cite{lewisbio1979,loc40,gnd40}。按他后来的自述，两个问题困扰他已久\cite{lewisbio1979}。一个是钢与咖啡的相对价格由什么决定。另一个是英国工业革命的头五十年里，实际工资大致不变，利润与储蓄却猛增\cite{lewisbio1979}。那天他忽然想到，两个问题有同一个答案。扔掉劳动数量固定的新古典假定，“无限的劳动供给”会压住工资，在前一个问题里带来便宜的咖啡，在后一个问题里带来高额的利润\cite{lewisbio1979}。
% 据: lb_bangkok, lb_steel, lb_britain, lb_unlimited

1954 年 5 月，《曼彻斯特学派》刊出他的文章《劳动无限供给条件下的经济发展》\cite{lewis1954}。他说，文章按古典的传统写成，作古典的假定，问古典的问题\cite{lewis1954}。按他的说法，经济发展理论的中心问题，是理解储蓄怎样相对于国民收入提高起来\cite{lewis1954}。一个社会过去只把国民收入的 4\% 或 5\% 甚至更少用于储蓄和投资。后来，它的自愿储蓄达到国民收入的 12\% 到 15\% 甚至更多\cite{lewis1954}。他说，文章发表以后，掌声与愤怒的叫声一样多\cite{lewisbio1979}。
% 据: lw_issue, lw_classical, lw_central, lw_why, lb_applause

上一章的索洛模型回答的是，一个穷国如果像富国那样储蓄和投资，会走到哪里（\chref{ch:solow}）。它把储蓄率当作给定。刘易斯要解释的正是储蓄率本身怎样提高（本书的解读）。1979 年，他与舒尔茨（Theodore W. Schultz，1902—1998）各得一半诺贝尔经济学奖\cite{nobel1979,loc41,gnd41}。获奖的理由是他们在经济发展上的开创性研究，特别考虑了发展中国家的问题\cite{nobel1979}。穷国怎样起飞？
% 据: nf_share, nf_motivation, loc_schultz, gnd_schultz

\section{当时的世界}\label{sec:development-world}

刘易斯生在圣卢西亚，1948 年 33 岁时成了曼彻斯特大学的正教授\cite{lewisbio1979}。按他的自述，曼彻斯特大批亚洲与非洲的学生，使他约从 1950 年起系统地讲授发展经济学\cite{lewisbio1979}。他的听众来自正在独立或刚刚独立的国家。
% 据: lb_born, lb_professor, lb_students

两次大战之间，出口初级产品的国家先吃了苦头。按世界银行《1987 年世界发展报告》的回顾，那些年它们的出口市场与进口工业品的渠道都缩小了，贸易条件急剧恶化\cite{wdr1987}。阿根廷等国因此在 1930 年代走上保护性的进口替代（\chref{ch:depressionworld}）\cite{wdr1987}。报告说，战后头二十年，初级产品贸易条件长期恶化的预测很普遍，许多新独立的国家于是想尽快转向工业\cite{wdr1987}。
% 据: wd_wars, wd_argentina, wd_secular, wd_secular2, wd_newly

报告还列出别的理由\cite{wdr1987}。许多新国家把专门生产初级产品同殖民地的身份连在一起，独立就要脱离殖民地的经济体系。许多人相信，欧洲、北美与日本的早期工业化靠过进口限制。苏联的快速工业化又让一些国家把进口替代同国有、计划与政府生产结合起来\cite{wdr1987}。报告说，这些国家的政府还限制国内贸易，欧洲与日本的早期工业化却得益于国内贸易限制的放松\cite{wdr1987}。苏联工业化的代价与那场饥荒在\chref{ch:soviet}。
% 据: wd_colonial, wd_protection, wd_soviet, wd_domestic

储蓄从哪里来，是另一层背景。按 Maddison 的估计，殖民统治把印度四分之一左右的储蓄转出了国（\chref{ch:periphery}）\cite{maddison1971}。刘易斯看的是国内的分配。他说，不发达经济的储蓄相对于国民收入低，原因不在人民穷，在资本家的利润相对于国民收入低\cite{lewis1954}。他还说，几乎所有的储蓄都来自得到利润或租金的人，工人的储蓄很少\cite{lewis1954}。要问的只是收入最高的那 10\% 的人，在剩余劳动的国家，他们拿到国民收入的 40\% 之多\cite{lewis1954}。在较发达的国家，这个比例接近 30\%。
% 据: md_quarter, lw_poor, lw_savers, lw_top10

刘易斯 1979 年回顾，1940 与 1950 年代的经济学家造了一整套理论，平衡增长、区域一体化、两缺口模型、结构性通货膨胀，都是在世界贸易停滞时才说得通\cite{lewis1979}。许多国家的政策也建立在这个假定上，向内，主要是进口替代\cite{lewis1979}。他说，1913—1939 年世界贸易每年只增长 0.9\%\cite{lewis1979}。按他的说法，贸易其实已经增长得很快，这一点要到 1960 年代后半期才被普遍承认\cite{lewis1979}。
% 据: ll_theories, ll_inwards, ll_trade, ll_recognised

发展也有了数字的目标。按刘易斯的说法，联合国在 1950 年代末把每年增长 5\% 定为 1960 年代的目标，所谓目标，是做不到却能鼓舞人的东西\cite{lewis1979}。出乎大家的意料，联合国的数字显示，到 1960 年代中期，平均增长率已经是 5\%\cite{lewis1979}。1970 年代的目标提到 6\%，1970 年代初又已经达到\cite{lewis1979}。联合国正准备把 1980 年代的目标定为 7\%，1974 年的衰退开始了\cite{lewis1979}。
% 据: ll_target60, ll_surprise, ll_target70, ll_target80

刘易斯 1979 年说，刚过去的二十年里，世界贸易每年实际增长约 8\%，速度是以前最快时的两倍\cite{lewis1979}。1873—1913 年，它每年的增长不到 4\%\cite{lewis1979}。
% 据: ll_eight, ll_1873

“起飞”是罗斯托（W. W. Rostow，1916—2003）的说法\cite{loc41,gnd41}。按 Easterly 的转述，罗斯托 1960 年的《经济增长的阶段》把起飞的标志定为投资从收入的 5\% 提到 10\%，几乎就是刘易斯六年前的说法\cite{easterly1997}。按 Easterly 的说法，罗斯托用哈罗德—多马模型算起飞需要的投资，再让西方的援助填上它与国民储蓄之间的“资金缺口”\cite{easterly1997}。这个模型在\chref{ch:solow}，它的检验在\secref{sec:development-push}。
% 据: loc_rostow, gnd_rostow, ea_stages, ea_takeoff, ea_rostow_gap

\section{二元经济：刘易斯的模型}\label{sec:development-lewis}

码头上等活的工人，抢着替人提包的年轻人，打零工的园丁，人数常常是活计所需的好几倍\cite{lewis1954}。刘易斯用他们说明，劳动过剩不只在农村。按他的定义，人口相对于资本与自然资源多到一定程度，经济里就有大片部门的劳动边际生产率可以忽略、为零甚至为负，劳动的供给是无限的\cite{lewis1954}。
% 据: lw_casual, lw_unlimited

\subsubsection{两个部门}

刘易斯把经济分成两个部门\cite{lewis1954}。资本主义部门使用可再生产的资本，并为此向资本家付报酬。维持生计的部门是其余的部分，不使用可再生产的资本，人均产出因此较低\cite{lewis1954}。他说，边际生产率是零还是很小，对他的分析并不根本。要紧的是工资定在维持生计的水平上，按这个价格，劳动的供给超过需求\cite{lewis1954}。
% 据: lw_capitalist, lw_subsistence, lw_notfundamental

维持生计部门的收入给资本主义部门的工资定了下限\cite{lewis1954}。他说，实际上两者之间通常差 30\% 或更多，一部分是城里的生活更贵，一部分是离开轻松的生活方式、进入更有纪律的城市环境要有补偿\cite{lewis1954}。资本主义部门雇人，雇到劳动的边际生产率等于工资为止，其余的产出是资本家的剩余\cite{lewis1954}。
% 据: lw_floor, lw_gap, lw_explain, lw_margin

在多数人是在自己土地上耕作的农民的经济里，劳动的最低价格有一个更客观的标准，就是农民的平均产品\cite{lewis1954}。农民在外面挣的如果不如留在家里能消费的多，就不会离开家庭农场。他也说，妇女的工作从家庭转到商业性就业，是经济发展最显著的特征之一\cite{lewis1954}。
% 据: lw_farmer, lw_farm, lw_women

用一个教学的形式写出来。资本主义部门的产出是$Y_c=A\,K^{\alpha}L_c^{1-\alpha}$，工资固定为$\bar w$，于是
\begin{keyeq}
(1-\alpha)\,A\,K^{\alpha}L_c^{-\alpha}=\bar w ,\qquad \Pi=Y_c-\bar w L_c=\alpha\,Y_c ,\qquad \frac{S}{Y}=\frac{\alpha\,Y_c}{Y_c+Y_s} .\label{eq:development-lewis}
\end{keyeq}
第一式是雇用的条件，工资不变时，雇用的人数$L_c$与资本$K$同比例增加。第二式是利润，第三式假定利润全部储蓄，$Y_s$是维持生计部门的产出（本书的推导）。刘易斯本人用的是图，没有写这个函数。
% 据: lw_margin

刘易斯说，剩余再投资，资本主义部门扩大，从维持生计部门吸收更多的人\cite{lewis1954}。剩余更大，资本形成也更多，这个过程一直继续到剩余劳动消失为止\cite{lewis1954}。在他的总结里，资本形成与技术进步提高的是利润在国民收入中的份额，工资没有提高\cite{lewis1954}。资本主义部门扩大，利润相对增长，国民收入中再投资的比例越来越大\cite{lewis1954}。按他的说法，经济发展的中心事实，是收入分配向储蓄的阶级倾斜\cite{lewis1954}。
% 据: lw_process, lw_continue, lw_share, lw_grow, lw_fact

在这个模型里，工资被压在维持生计的水平上，增长的好处先归利润。刘易斯也直说，资本主义部门的工资取决于维持生计部门的收入，这有时有巨大的政治意义，因为资本家有直接的利益去压低维持生计部门工人的生产率\cite{lewis1954}。资本积累赶上劳动供给以后，工资开始高于维持生计的水平，他的封闭经济模型就不再成立\cite{lewis1954}。本书把这个时点叫作转折点。转折点以后，这个国家仍被别的剩余劳动国家包围。他说，工资一开始上升，大规模移民与资本输出就会抑制它\cite{lewis1954}。
% 据: lw_political, lw_catchup, lw_noclosed, lw_immigration

资本也可以来自信用创造。他说，在这个模型里，通货膨胀创造的资本实际成本为零，与用利润创造的资本一样有用\cite{lewis1954}。为创造生产性资本的通货膨胀会自我消解，价格随资本形成上涨，产出上市后又回落\cite{lewis1954}。为战争筹集资源的通货膨胀却可能累积（\chref{ch:hyperinflation}）。
% 据: lw_credit, lw_inflation

刘易斯把同样的逻辑用到世界经济上。他在总结里说，热带的商业产品按它所提供的生活水平算这么便宜，主要原因是热带粮食生产的人均效率低\cite{lewis1954}。出口产业效率提高的好处，几乎全部归了外国的消费者\cite{lewis1954}。这与下一节普雷维什的结论相近，两人的机制却不同（本书的解读）。
% 据: lw_tropical, lw_foreign

\subsubsection{储蓄率怎样上升}

刘易斯的图就是这样画的\cite{lewis1954}。横轴是劳动的数量，OS 是维持生计的平均收入，OW 是资本主义部门的工资。剩余再投资，固定资本增加，劳动的边际生产率曲线整个抬高，雇用与剩余都变大\cite{lewis1954}。
% 据: lw_diagram, lw_raised

\exref{ex:development-lewis}用教学输入算这个过程（\figref{fig:development_lewis}）。资本主义部门开始只雇 8\% 的劳动力，全社会的储蓄率不到 5\%。工资不变，利润全部再投资，资本主义部门每年扩大约 11\%，约 10 年后吸收完剩余劳动，储蓄率升到约 13\%（本书的计算）。国民收入增加的部分，三分之二以上是利润，工人的收入只因有人换到工资更高的部门才增加（本书的计算）。

\begin{example}{刘易斯的转折点}\label{ex:development-lewis}
教学输入：劳动力 100 人，维持生计部门每人生产并挣得 1，资本主义部门的工资$\bar w=1.3$（刘易斯说的差距 30\%），$\alpha=1/3$，资本主义部门的资本—产出比为 3，利润全部再投资，没有折旧，人口不变。资本主义部门开始雇 8 人，能按这个工资招到的人以 24 人为限。（1）开始时全社会的储蓄率是多少？（2）第几年到达转折点？（3）那时的储蓄率是多少，国民收入增加的部分里利润占多少？
\solution
（1）$Y_c=\bar w L_c/(1-\alpha)=1.95\times8=15.6$，利润$5.2$，$Y=15.6+92=107.6$，储蓄率约 4.8\%。（2）资本每年增加$\alpha/3=1/9$，雇用的人数同比例增加。从 8 人到 24 人，要$\ln3/\ln(10/9)\approx10.4$ 年，即第 11 年到达。（3）$Y_c=46.8$，利润$15.6$，$Y=46.8+76=122.8$，储蓄率约 12.7\%。国民收入增加 15.2，其中利润增加 10.4，约占 68\%；换到资本主义部门的 16 人每人多挣 0.3，共 4.8（本书的计算）。
\end{example}

\fig{development_lewis}{刘易斯模型（本书的计算，教学输入同\exref{ex:development-lewis}）。(a) 资本主义部门的劳动边际产品曲线随资本积累右移，工资线是平的，交点左边、工资线以上的面积是利润。(b) 储蓄率与资本主义部门的就业份额，竖线是转折点。}{fig:development_lewis}

\begin{pitfall}[劳动过剩就是农业劳动的边际产品为零]
刘易斯自己说，边际生产率是零还是很小，对他的分析并不根本\cite{lewis1954}。模型需要的只是，按一个固定的工资招得到的人比需要的多。他说这样的人不只在农村，小买卖与家务劳动里也有\cite{lewis1954}。舒尔茨 1979 年的演讲说，农民在自己小小的配置范围里是精于微调的企业家，许多专家看不出他们多有效率\cite{schultz1979}。农民对价格敏感，与工资被维持生计的收入压住，两件事可以同时成立（本书的解读）。
% 据: lw_notfundamental, lw_petty, lw_domestic, sc_finetuning
\end{pitfall}

\section{贸易条件与进口替代}\label{sec:development-isi}

1949 年 10 月 9 日，联合国拉丁美洲经济委员会的执行秘书在智利圣地亚哥写信给联合国秘书长，转交一篇论文\cite{prebisch1950}。论文的作者是普雷维什（Raúl Prebisch，1901—1986）\cite{prebisch1950,bnf29}，他在\chref{ch:depressionworld}里做过阿根廷中央银行的总经理。1950 年 4 月 27 日，联合国在纽约成功湖印出它的英文本，原文是西班牙文\cite{prebisch1950}。按送达信的说法，经济及社会理事会设立这个委员会时，把发起或主持拉美经济问题的调查研究定为它的主要目的之一\cite{prebisch1950}。
% 据: pb_letter, pb_essay, pb_author, pb_secretary, pb_symbol, pb_date, pb_spanish, pb_lake, pb_purpose, bnf_prebisch

报告开头说，拉丁美洲的现实正在瓦解过时的国际分工图式\cite{prebisch1950}。按这个图式，拉丁美洲作为世界经济体系的外围，任务是为大工业中心生产粮食和原料。报告说，这个图式假定技术进步的好处会通过降价或提高收入分到整个共同体，事实证明这个假定是错的\cite{prebisch1950}。生产率提高的好处，没有以可比的程度到达外围。报告说，工业化本身“不是目的”，是这些国家分享技术进步的好处、逐步提高大众生活水平的主要手段\cite{prebisch1950}。
% 据: pb_schema, pb_periphery, pb_assumption, pb_false, pb_benefits, pb_means

\subsubsection{普雷维什的数}

报告的表 1 用英国贸易部的进出口平均价格指数，以 1876—1880 年为 100，算初级产品与工业品的价格之比\cite{prebisch1950}。这个比到 1936—1938 年降到 64.1，1946—1947 年回到 68.7（\figref{fig:development_terms}）\cite{prebisch1950}。报告说，同样多的初级产品，1930 年代只能买到 1860 年代工业品的 63\%，平均要多 58.6\% 的初级产品才买得到同样多的工业品\cite{prebisch1950}。它也承认，价格指数没有反映工业品质量的差别\cite{prebisch1950}。报告第 12 页又说，从 1870 年代到 1930 年代，价格之比恶化了 36.5\%\cite{prebisch1950}。表里的比值在 1880 年代前半期最高，为 102.4，此后一路走低。1921—1925 年是 67.3。1931—1935 年最低，为 62.0\cite{prebisch1950}。
% 据: pb_t1_title, pb_t1_source, pb_t1_all, pb_63, pb_quality, pb_365

出口价格指数$P_x$与进口价格指数$P_m$之比$P_x/P_m$，叫作净易货贸易条件。再乘上出口量$Q_x$，$(P_x/P_m)\,Q_x$叫作收入贸易条件，是出口能换到的进口量（本书的概括）。价格之比下降，出口量增长得够快，能买到的进口仍可以增加。普雷维什的表只看价格之比。

报告把原因放在周期里\cite{prebisch1950}。上升时，中心的一部分利润被工资吸收，这来自企业家之间的竞争与工会的压力。下降时，中心的工资因为人所共知的抵制降不下来，压力就转向外围\cite{prebisch1950}。报告说，初级生产的工人特有的缺乏组织，使他们得不到与工业国可比的工资增长\cite{prebisch1950}。于是初级产品的价格在上升时涨得多，在下降时跌得更多，差距在一个个周期里逐步拉大\cite{prebisch1950}。
% 据: pb_upswing, pb_rigid, pb_organization, pb_cycles

\fig{development_terms}{普雷维什报告表 1：初级产品与工业品的价格之比（1876—1880 年＝100）。数据是英国贸易部的进出口平均价格指数，1930 年代只到 1938 年。横轴按各时期的中点，两次大战期间缺数。}{fig:development_terms}

\subsubsection{有效保护率}

拉丁美洲、印度与许多新国家用关税、进口许可与高估的汇率保护本国工业。世界银行 1987 年的报告说，进口替代的国家通常先生产最终工业品来替代进口，许多国家起初制造业增长很快\cite{wdr1987}。生产要用进口的中间品与资本品，持续的工业增长靠出口挣外汇\cite{wdr1987}。
% 据: wd_final, wd_foreign

按世界银行的报告，内向型战略的工具不只是关税\cite{wdr1987}。决策者往往更喜欢进口许可与数量限制这类直接控制。印花税、港口税与预付保证金这类隐蔽的进口税也很常见，国产化率的要求越来越多\cite{wdr1987}。许多发展中国家的公营企业扩张得很快，尤其在工业里\cite{wdr1987}。制造业的跨国公司被吸引到一些发展中大国，特别是拉丁美洲，原因之一是那里的贸易政策限制最终工业品的进口\cite{wdr1987}。
% 据: wd_tools, wd_hidden, wd_public, wd_multinationals

关税保护的是什么，要看增加值。报告把有效保护率定义为按国内价格算的单位增加值与按世界价格算的单位增加值之差，以后者的百分比表示\cite{wdr1987}。设最终产品的关税为$t_f$，投入品的关税为$t_m$，投入品按世界价格占产品价格的比例为$a$，
\begin{keyeq}
\text{ERP}=\frac{v'-v}{v}=\frac{t_f-a\,t_m}{1-a} ,\label{eq:development-erp}
\end{keyeq}
$v$与$v'$是按世界价格与国内价格算的单位增加值（本书的推导）。投入品的关税越低、增加值所占的比例越小，同样的名义关税给的有效保护越高。
% 据: wd_erp_def

\begin{example}{一件毛衣的有效保护}\label{ex:development-erp}
用世界银行报告的例子\cite{wdr1987}。没有进口限制时，一件毛衣卖 100 美元，羊毛与纽扣按世界价格值 60 美元，增加值 40 美元。（1）毛衣征 20\% 的关税，投入品免税，有效保护率是多少？（2）投入品也征 10\% 的关税呢？（3）投入品的关税提到 40\% 呢？
\solution
（1）国内的增加值是$120-60=60$，有效保护率$(60-40)/40=50\%$。（2）投入品值 66 美元，增加值 54，有效保护率 35\%，与式~\eqref{eq:development-erp}的$(0.2-0.6\times0.1)/0.4$相同\cite{wdr1987}。（3）$(0.2-0.6\times0.4)/0.4=-10\%$，保护变成了负的（本书的计算）。毛衣的买主每件多付 20 美元，毛衣厂的增加值却多出一半。
% 据: wd_sweater, wd_sweater60, wd_sweater20, wd_sweaterdom, wd_sweater50, wd_sweater35
\end{example}

关税与许可由政府定，账由谁付，世界银行报告的专栏有估计\cite{wdr1987}。它说，按能找到数据的国家，平均都有不利于出口的偏向，新加坡与韩国可能是例外\cite{wdr1987}。除韩国以外，每个国家都明显不利于初级产品或农业部门，偏向制造业。韩国的例外，是因为对稻米的价格支持不断提高\cite{wdr1987}。这种偏向在内向型的国家更显眼\cite{wdr1987}。以国内销售与出口的有效保护之差衡量，专栏里不利于出口的偏向最小的是哥伦比亚，1969 年为 9 个百分点\cite{wdr1987}。最大的是智利，1967 年达 229 个百分点。
% 据: wd_bias, wd_singapore, wd_agri, wd_inward, wd_range

舒尔茨 1979 年的演讲说，政策把最高的优先给了工业，包括让粮价保持便宜\cite{schultz1979}。按他的说法，原因是国内政治通常偏向城市人口，牺牲人数多得多的农村人口\cite{schultz1979}。关税向买进口品与国产品的消费者征税，高估的汇率向出口者征税，这里的出口者多是农民与矿业（本书的概括）。
% 据: sc_priority, sc_politics

\begin{compare}{进口替代与出口导向}{强内向型（进口替代）}{强外向型（出口导向）}
\cmprow{理由}{工业化是分享技术进步好处的主要手段，初级产品的价格之比长期下降\cite{prebisch1950}}{刘易斯 1979 年说，几乎每个国家后来都发现了出口的好处\cite{lewis1979}}
\cmprow{工具}{有效保护率高，直接控制与许可普遍，汇率明显高估\cite{wdr1987}}{贸易控制很少或很低，进口品与出口品的有效汇率大致相等\cite{wdr1987}}
\cmprow{1963—1985 年的成员}{阿根廷、加纳、印度等\cite{wdr1987}}{香港、韩国、新加坡\cite{wdr1987}}
\cmprow{1963—1973 年每年的人均 GNP 增长}{1.6\%\cite{wdr1987}}{6.9\%\cite{wdr1987}}
\cmprow{1973—1985 年}{平均每年下降 0.1\%\cite{wdr1987}}{5.9\%\cite{wdr1987}}
\cmpsum{报告自己说，贸易战略与宏观表现的联系并不完全清楚，是外向带来了更好的表现，还是更好的表现为外向铺了路\cite{wdr1987}。东亚的证据在\chref{ch:eastasia}。}
% 据: wd_inward_def, wd_controls, wd_overvalued, wd_outward_def, wd_parity, wd_groups, wd_gnp63, wd_gnp63b, wd_gnp73a, wd_gnp73b, wd_unclear
\end{compare}

\section{大推进与资金缺口}\label{sec:development-push}

一家工厂单独开张，可能亏本，许多工厂一起开张，却可能都赚钱。罗森斯坦—罗丹（Paul N. Rosenstein-Rodan，1902—1985）1943 年提出过这个想法\cite{msv1988,loc41,gnd41}。按 Murphy、Shleifer 与 Vishny 1988 年的工作论文的转述，各部门同时采用规模报酬递增的技术，各自创造的收入就成为别的部门的需求，市场随之扩大\cite{msv1988}。即使没有一个部门单独工业化能够保本，许多部门同时工业化也可以自我维持。这个想法后来发展成平衡增长或大推进的学说\cite{msv1988}。
% 据: ms_rr, ms_rr2, ms_simultaneous, ms_doctrine, loc_rr, gnd_rr

\subsubsection{两个均衡}

他们的模型里，每种商品可以用手工生产，一单位劳动出一单位产品\cite{msv1988}。工厂先要投入$F$单位劳动的固定成本，此后一单位劳动出$a>1$单位产品，工资要比手工多出补偿$v$\cite{msv1988}。总收入为$y$时，一家工厂的利润是$y\,(1-(1+v)/a)-F\,(1+v)$。设劳动力为$L$，两个均衡同时存在的条件是
\begin{keyeq}
L\Bigl(1-\frac{1+v}{a}\Bigr)-F\,(1+v)<0 ,\qquad a\,(L-F)-L\,(1+v)>0 .\label{eq:development-push}
\end{keyeq}
第一式说，大家都不建厂时，一家单独建厂会亏本。第二式说，大家都建厂时，每家都赚钱\cite{msv1988}。他们用$\alpha$表示生产率，这里写作$a$。模型要有意思，生产率的提高必须超过补偿，即$a-1>v$，否则工厂把全部效率收益都给工人也招不到人\cite{msv1988}。
% 据: ms_cottage, ms_irs, ms_wage, ms_profit, ms_eq9, ms_eq10, ms_eq8

为什么一家亏本，大家却都赚钱？按他们的解释，建厂的企业付了工资溢价，即使自己亏本，也扩大了别的工业品的市场\cite{msv1988}。企业的利润没有把多付的工资算进它对总需求的贡献。工业化的均衡对所有人都更好，工人得到同样的效用，还多了利润\cite{msv1988}。协调各部门的投资，于是成了政府可以推动的事（\exref{ex:development-push}）。
% 据: ms_premium, ms_pareto, ms_coordination

论文还给了两条别的渠道\cite{msv1988}。今天投资、明天低成本生产的企业，把收入与需求移到将来，使别的企业也更愿意今天投资。用铁路这类共用基础设施的企业，帮别的部门分摊修路的固定成本\cite{msv1988}。按他们的说法，这些传递机制对欠发达国家似乎都有一定的意义\cite{msv1988}。
% 据: ms_dynamic, ms_railroad, ms_relevance

在印度，协调各部门投资的是计划委员会。1958 年的摘要说，国家发展委员会要求在分配项目时照顾欠发达地区的需要，并给社会服务与社区发展以很高的优先\cite{india1958}。
% 据: i8_areas

\begin{example}{两个均衡}\label{ex:development-push}
教学输入：$L=100$，$a=1.5$，$v=0.2$。（1）$F$在什么范围内，式~\eqref{eq:development-push}的两个条件同时成立？（2）取$F=18$，没人建厂时单独建厂的利润是多少？大家都建厂时每家的利润是多少？
\solution
（1）第一式要$100\times(1-1.2/1.5)=20<1.2F$，即$F>16.7$；第二式要$1.5\times(100-F)>120$，即$F<20$。（2）单独建厂：$20-1.2\times18=-1.6$，亏本。大家都建厂：$1.5\times82-120=3$，每家都赚（本书的计算）。
\end{example}

\subsubsection{资金缺口}

计划者算投资，常用的是哈罗德—多马的一条规则（\chref{ch:solow}）。目标增长率$g$乘上增量资本产出比（ICOR），就是需要的投资率，再减去国内储蓄率，就是要从外面借来的缺口：
\begin{keyeq}
\frac{I}{Y}=\text{ICOR}\times g ,\qquad \text{缺口}=\frac{I}{Y}-\frac{S}{Y} .\label{eq:development-gap}
\end{keyeq}
按 Easterly 的转述，当时的经济学家认为 ICOR 在 2 到 5 之间\cite{easterly1997}。他的例子是，投资率 4\%、ICOR 为 4 时，增长率每年只有 1\%\cite{easterly1997}。这赶不上 2\% 的人口增长。投资提到刘易斯说的 12\%，增长率是 3\%\cite{easterly1997}。国民储蓄仍是 4\%，差的 8 个百分点由西方的援助去填\cite{easterly1997}。
% 据: ea_icor, ea_example, ea_gap12, ea_gap

这条规则后来进了国际机构。按 Easterly 的转述，Chenery 与 Strout 1966 年的著名应用更强调国民储蓄\cite{easterly1997}。1971 年，世界银行把 Chenery 版本的哈罗德—多马模型编成了计算机程序，当时 Chenery 是行长麦克纳马拉的首席经济顾问\cite{easterly1997}。
% 据: ea_chenery, ea_rmsm

\begin{example}{印度第三个五年计划的缺口}\label{ex:development-gap}
印度计划委员会 1961 年说，要每年增长 5\% 以上，净投资要超过国民收入的 14\%，当时约是 11.5\%\cite{india1961}。国内储蓄要从约 8.5\% 提到计划末的约 11.5\%，其余靠外援\cite{india1961}。（1）计划隐含的 ICOR 是多少？（2）计划末的缺口是多少？（3）ICOR 若是世界银行报告里强内向型国家 1973—1985 年的 8.7，14\% 的投资能带来多少增长？
\solution
（1）$14/5=2.8$。（2）$14-11.5=2.5$个百分点，约为国民收入的 2.5\%。（3）$14/8.7\approx1.6\%$（本书的计算）\cite{wdr1987}。按\chref{ch:solow}的模型，ICOR 并不是常数，资本越多，同样的投资带来的增长越少。
% 据: i3_growth, i3_savings, wd_icor73
\end{example}

按 Easterly 的检验，88 个国家 1965—1995 年的数据里，援助与投资的关系 60\% 是负的\cite{easterly1997}。只有 6 个国家通过检验，即援助增加投资至少同样多，而且显著\cite{easterly1997}。整个资金缺口模型只适合突尼斯一国。按他的说法，138 国里成功 1 个，很可能只是偶然\cite{easterly1997}。按他的算法，赞比亚的援助如果都变成投资、又按 ICOR 带来增长，它今天会是工业国，而它是世界上最穷的国家之一\cite{easterly1997}。
% 据: ea_88, ea_60, ea_six, ea_tunisia, ea_zambia

乐观的预测也落空过。按 Easterly 的转述，Kamarck 1967 年预测，加纳、尼日利亚与赞比亚等矿产给了它们储蓄的国家，增长率很可能达到或超过 7\%\cite{easterly1997}。这些国家 1967—1992 年增长率的中位数只有 2.8\%\cite{easterly1997}。按他的表 4，把增长 7\% 以上的年份挑出来，ICOR 取通常的 3.5 时，九成以上没有达到所谓必要的投资率\cite{easterly1997}。
% 据: ea_mineral, ea_kamarck, ea_median, ea_table4

\section{印度与加纳：谁决定，谁承担}\label{sec:development-cases}

1955 年 3 月，统计学家马哈拉诺比斯（P. C. Mahalanobis，1893—1972）提出第二个五年计划的建议草案，通称“计划框架”\cite{india1956,loc41,gnd41}。按计划委员会 1956 年 2 月的草案大纲，国家发展委员会 5 月初大体同意它的思路\cite{india1956}。委员会要求计划能在五年里让国民收入增加约 25\%，给 1,000 万到 1,200 万人提供就业，并体现社会主义类型社会的政策决定\cite{india1956}。按大纲的数，人均收入要从 1955—1956 年度的 280 卢比提到 1960—1961 年度的 330 卢比，卢比是印度的货币\cite{india1956}。计划的目标之一是快速工业化，特别强调“基础与重工业”\cite{india1956}。
% 据: in_planframe, in_planframe2, in_date, in_ndc, in_agreed, in_targets, in_socialistic, in_percapita, in_heavy, loc_mahalanobis, gnd_mahalanobis

计划委员会知道这条路的代价\cite{india1956}。大纲说，印度在钢这样的部门可能有比较成本优势，应当尽量发展重工业与资本品工业。它也说，这样的工业发展需要大量资本，吸收的人力却比较少\cite{india1956}。投资于基础工业会增加对消费品的需求，短期里却不增加消费品的供给。大纲说，这表明在从基础加强经济的时候，消费上作出一定的牺牲不可避免\cite{india1956}。
% 据: in_steel, in_capital, in_demand, in_sacrifice

投资率也由计划者定。大纲要把投资从国民收入的约 7\% 提到第二个计划末的近 12\%\cite{india1956}。国内储蓄要从约 7\% 提到 10\% 以上，其余靠外部资源\cite{india1956}。工资低的国家为什么要用为高工资设计的技术，先建资本密集的重工业？这是\chref{ch:divergence}末尾反过来问的 Allen 的问题。计划者的回答是长远的产能，代价由眼前的消费与没被吸收的人力承担（本书的解读）。按世界银行 1987 年的分类，印度 1963—1985 年属于强内向型\cite{wdr1987}。
% 据: in_investment, in_savings, wd_groups

外汇先成了瓶颈。按计划委员会 1958 年 9 月的摘要，它在当年 5 月向国家发展委员会和议会提交备忘录，把计划分成两部分\cite{india1958}。A 部分除了与农业增产直接相关的项目，还有“核心”项目\cite{india1958}。它的总支出限于 4,500 克罗尔卢比\cite{india1958}。克罗尔是印度的计数单位，1 克罗尔等于 1,000 万\cite{india1956}。外汇储备的下降，在 1958 年 3 月提交议会的一份文件里已有详细检讨\cite{india1958}。
% 据: i8_date, i8_memo, i8_parliament, i8_parts, i8_core, i8_ceiling, in_crore, i8_reserves

1961 年的第三个计划说，国家要在五年里做到前两个计划十年里做到的那么多\cite{india1961}。剩余劳动却没有被重工业吸收。这个计划估计，计划里的项目可以提供约 1,400 万个新增就业，同期劳动力却要增加约 1,700 万\cite{india1961}。不足的部分，计划打算靠农村的大规模工程、乡村与小工业等办法解决\cite{india1961}。
% 据: i3_ambition, i3_employment, i3_balance

加纳 1957 年独立，恩克鲁玛（Kwame Nkrumah，1909—1972）领导它，直到 1966 年被政变推翻\cite{loc41,gnd41}。1957—1963 年，刘易斯先后做过加纳总理的联合国经济顾问、联合国特别基金的副总干事与西印度大学的校长\cite{lewisbio1979}。恩克鲁玛 1957 年 3 月出任总理，1960 年加纳成为共和国，他改任总统\cite{loc41}。1963 年 12 月，计划委员会办公室印出 1963/64—1969/70 年七年发展计划的简要大纲，前言署名的是恩克鲁玛总统\cite{ghana1963}。前言说，计划要用科学技术使农业与工业发生革命\cite{ghana1963}。
% 据: loc_nkrumah, gnd_nkrumah, gnd_nkrumah_d, lb_ghana, loc_nkrumah_pm, gh_title, gh_date, gh_foreword, gh_science

大纲说，十个就业的加纳人里有六个在农业，农业生产了一半以上的国民收入\cite{ghana1963}。劳动从土地转向工厂的时候，农业部门要为工业品提供国内市场，还要提供工业发展所需储蓄的相当大一部分\cite{ghana1963}。它说，经济发展成败的重担落在从事农业的人身上，农业的这种转变是国民收入持续增长的先决条件\cite{ghana1963}。七年的总投资估计为 10.16 亿加纳镑，其中政府 4.76 亿，私人 5.40 亿\cite{ghana1963}。加纳镑（£G）是 1965 年 7 月以前的加纳货币\cite{wbstab1967}。政府的部分有 8,700 万靠短期借款，即赤字财政\cite{ghana1963}。可可价格按每吨 190 加纳镑的平均数估计，大纲说，可可涨价会增加政府收入，减少赤字财政，赤字财政只是余项\cite{ghana1963}。
% 据: gh_six, gh_savings, gh_burden, gh_prerequisite, gh_total, gh_split, ws_pound, gh_deficit, gh_cocoa, gh_residual

\begin{case}{加纳 · 1947—1970 · 可可销售局}
付账的是种可可的农民。按世界银行 1970 年的项目报告，可可销售局按 1947 年的条例设立，可以定收购价、积累盈余，在世界价格低时用盈余支持收购价\cite{wbcocoa1970}。按报告用新塞地计的数，它的价格稳定储备 1957 年约有 9,000 万，这笔钱和后来积累的资金都被政府用于可可以外的用途\cite{wbcocoa1970}。新塞地（NC）是 1967 年 2 月起的加纳货币\cite{wbstab1967}。1965 年的新法律禁止销售局积累储备，要它把经营盈余交给政府\cite{wbcocoa1970}。
% 据: wc_cmb, wc_support, wc_reserve, wc_noncocoa, ws_newcedi, wc_1965

报告说，生产者价格 1959 年是每磅 12 美分，1 磅约合 0.45 公斤\cite{wbcocoa1970}。到 1965 年，价格降到 6.7 美分，新种植减少，农场疏于管理\cite{wbcocoa1970}。可可农药的用量 1962 年最多，有 31.2 万加仑，1966 年只剩 2.7 万加仑\cite{wbcocoa1970}。收成在 1964/65 年度达到 568,032 吨，此后下降，近年平均约 40 万吨\cite{wbcocoa1970}。1966 年政府更迭以后，收购价提到每磅 13 美分\cite{wbcocoa1970}。
% 据: wc_1959, wc_1965price, wc_plantings, wc_pesticide, wc_crop, wc_400, wc_change, wc_13, wc_lb

世界银行按每磅 27 美分的长期价格算，农民得 13 美分，销售局的盈余 1.5 美分，出口税 7.7 美分\cite{wbcocoa1970}。收成若是 40 万吨，政府每年留下约 8,000 万新塞地。报告说，留下这么大的份额，只有在加纳眼下困难的经济与财政状况下才说得过去\cite{wbcocoa1970}。按世界价格卖可可的是国家，拿差价去投资的也是国家，少收的是农民（本书的解读）。
% 据: wc_split, wc_duty, wc_80, wc_justifiable
\end{case}

两国的账可以并排看。印度的计划委员会让消费为重工业让路，加纳的计划要农业提供国内市场和储蓄。

\begin{compare}{两个计划}{印度的第二个五年计划}{加纳的七年发展计划}
\cmprow{文件}{计划委员会的草案大纲，1956 年 2 月\cite{india1956}}{计划委员会办公室的简要大纲，1963 年 12 月\cite{ghana1963}}
\cmprow{拍板的}{国家发展委员会大体同意，修订再提交它与议会\cite{india1956,india1958}}{前言署名的恩克鲁玛总统\cite{ghana1963}}
\cmprow{优先}{基础与重工业\cite{india1956}}{用科学技术使农业与工业发生革命\cite{ghana1963}}
\cmprow{钱从哪里来}{提高国内储蓄，其余靠外部资源\cite{india1956}}{农业提供储蓄的相当大一部分，政府的部分有一块靠赤字财政\cite{ghana1963}}
\cmprow{谁承担}{眼前的消费与没被吸收的人力（本书的解读）}{种可可的农民（本书的解读）}
\cmpsum{决定投资率与优先次序的是计划委员会与政府，付账的是消费者与种可可的农民（本书的概括）。按世界银行 1987 年的分类，加纳与印度 1963—1985 年都属于强内向型\cite{wdr1987}。}
% 据: in_date, in_agreed, i8_memo, i8_parliament, gh_date, gh_foreword, in_heavy, gh_science, in_savings, gh_savings, gh_deficit, wd_groups
\end{compare}

\section{数据检验：起飞了吗}\label{sec:development-check}

1960 年，韩国的人均 GDP 约为加纳的七成，1980 年已是加纳的三倍多（\tabref{tab:development-paths}，本书的计算）\cite{mpd2023,mpd41src}。这二十年里，韩国每年增长约 7.1\%，加纳每年下降约 0.9\%（本书的计算）。世界银行把韩国 1963—1985 年归入强外向型，把加纳归入强内向型\cite{wdr1987}。两国的起点相近，贸易政策的方向相反，结果相差很远。单凭这一对比还不能断定原因，两国在起点以外的差别还很多（本书的解读）。国际元是按 2011 年购买力平价折算的记账单位，即按同一篮子商品在各国的价格比较\cite{mpd2023}。1950 年，阿根廷的人均 GDP 约是韩国的八倍，1990 年韩国已经超过阿根廷（本书的计算）。上一章的检验里，1960 年最穷的四分之一共 27 个国家，有 13 个此后增长得比最富的四分之一的平均速度慢（\figref{fig:solow_convergence}）\cite{pwt1001}。
% 据: wd_groups

\begin{table}[tbp]
\centering\small
\caption{六个经济体的人均 GDP（2011 年国际元）与每年的增长率（本书的计算）}\label{tab:development-paths}
\begin{tabularx}{\linewidth}{@{}Y r r r r r@{}}
\toprule
 & 1950 & 1980 & 2000 & 1950—1980（\%） & 1980—2000（\%） \\
\midrule
阿根廷 & 7,949 & 13,080 & 14,369 & 1.7 & 0.5 \\
巴西 & 2,236 & 8,249 & 9,834 & 4.4 & 0.9 \\
墨西哥 & 3,510 & 9,929 & 12,613 & 3.5 & 1.2 \\
印度 & 987 & 1,495 & 2,753 & 1.4 & 3.1 \\
加纳 & 1,788 & 1,844 & 2,100 & 0.1 & 0.7 \\
韩国 & 998 & 6,064 & 23,108 & 6.2 & 6.9 \\
\bottomrule
\end{tabularx}
\datasource{Maddison Project Database 2023（Bolt 与 van Zanden 2024，CC BY 4.0），工作簿 mpd2023\_web.xlsx 的 Full data 表。各经济体的原始来源按数据库的引用要求列在参考文献里\cite{mpd41src}。增长率按复利计。}
\end{table}

数字不支持一个简单的结论。巴西与墨西哥在进口替代的年代里增长并不慢，1950—1980 年每年分别约 4.4\% 与 3.5\%（本书的计算）。世界银行把巴西 1963—1985 年归入中等外向型，把墨西哥归入中等内向型\cite{wdr1987}。账在后来才来，1980 年以后，三个拉美国家都增长得很慢（\figref{fig:development_paths}）。债务怎样为 1970 年代的进口替代付账，又怎样在 1982 年收场，在\chref{ch:latamdebt}。
% 据: wd_groups

印度与加纳的计划都把投资放在第一位，结果却不同。印度 1950—1980 年每年只增长约 1.4\%，1980 年以后加快。加纳的人均 GDP 在 1971 年以后下降\cite{mpd2023}。到 1983 年，它只有 1,487 国际元，比 1950 年还低（本书的计算）。刘易斯 1979 年说，到 1973 年为止的四分之一个世纪里，发展中国家与发达国家的 GDP 每年都增长约 5\%\cite{lewis1979}。发展中国家的人口增长更快，人均增长每年约 2.5\%，发达国家约 4\%\cite{lewis1979}。他说，按绝对水平看，发展中国家的表现很了不起，人均收入的差距却继续迅速拉大\cite{lewis1979}。
% 据: ll_five, ll_perhead, ll_remarkable

支持进口替代的一方也有数字。刘易斯 1979 年说，发展中国家的工业生产增长得比发达国家快\cite{lewis1979}。至少八个发展中国家的机器已占制造业产出的 15\% 以上，即印度、巴西、新加坡、智利、韩国、阿根廷、墨西哥与以色列\cite{lewis1979}。世界银行的数字显示，强外向型经济 1963 年的国内储蓄只占 GDP 的 13.0\%，是四组里最低的\cite{wdr1987}。到 1985 年，它升到 31.4\%。储蓄率并不是起飞以前就有的，它随增长一起上来（本书的解读）。
% 据: ll_industry, ll_machinery, wd_saving, wd_saving85

投资的效率也不同。按世界银行的数字，1963—1973 年外向型两组的增量资本产出比平均是 2.5\cite{wdr1987}。中等内向型与强内向型分别是 3.3 与 5.2。1973—1985 年各组都明显变差\cite{wdr1987}。外向型两组是 4.5 与 5.0。内向型两组是 6.2 与 8.7。同样的投资，在内向型经济里带来的增长少得多，这也是\exref{ex:development-gap}第（3）问的意思（本书的解读）。
% 据: wd_icor63, wd_icor63b, wd_icor73

\fig{development_paths}{六个经济体的人均 GDP，1950—2022 年（2011 年国际元，对数刻度）。(a) 三个拉美国家 1980 年以前增长，此后停滞。(b) 印度、加纳与韩国 1960 年的起点相近，此后分开。}{fig:development_paths}

\section{四个争论}\label{sec:development-debate}

\subsubsection{有没有剩余劳动}

第一个争论是剩余劳动。刘易斯说劳动的边际生产率在大片部门可以忽略，又说这一点并不根本\cite{lewis1954}。舒尔茨的回答是，农民是计算成本、收益与风险的经济人\cite{schultz1979}。按他的说法，最大的错误是以为标准的经济理论不足以理解低收入国家，为此造的模型一度广受称赞，后来才看清它们至多是知识上的“珍玩”\cite{schultz1979}。加纳的可可农民在收购价降低时少种、少喷农药，表现出的正是对价格的反应（本书的解读）。舒尔茨自己的保留是，农民的成就预示着有利的前景，前提是不被歧视农业的政治与政策耗掉\cite{schultz1979}。刘易斯说，二十五年里其他学者写了五本书和许多文章争论这个命题，争论还在继续\cite{lewisbio1979}。至今没有定论。
% 据: sc_calculating, sc_mistake, sc_curiosities, sc_prospects, lb_debate

\subsubsection{贸易条件是否长期恶化}

第二个争论是贸易条件。普雷维什的表从 1876—1880 年算起，用的是英国一国的进出口价格，1930 年代只算到 1938 年\cite{prebisch1950}。他自己说指数没有反映质量。世界银行 1987 年说，它第 2 章的数据也显示初级产品的贸易条件长期下降，但没有计入工业品质量的改进，要谨慎解读\cite{wdr1987}。报告说，初级产品的相对价格很可能在长期下降，问题是内向型的战略是不是对的回应\cite{wdr1987}。刘易斯 1979 年也说，贸易增长快时，贸易条件本应对发展中国家更有利，这一次却没有\cite{lewis1979}。他的推理给了另一种机制，发达国家增长放慢，发展中国家也跟着放慢，贸易条件就转向不利于后者\cite{lewis1979}。至今没有定论。
% 据: pb_t1_source, pb_t1_1938, wd_chapter2, wd_caution, wd_maywell, wd_response, ll_terms, ll_slower

\subsubsection{进口替代失败了吗}

第三个争论是进口替代本身。世界银行把 41 个经济体按贸易战略分组，外向型的几乎各方面都更好\cite{wdr1987}。同一份报告承认，分类的中间两组可能有争议，因果的方向也不清楚\cite{wdr1987}。报告说，早早转向出口扩张的国家，例如韩国，工业的增长持续了下去\cite{wdr1987}。刘易斯 1954 年的总结里说，比较成本的规律在剩余劳动的国家同样成立，那里它是主张保护的同样有效的理由\cite{lewis1954}。批评一方最有力的是分组与有效保护率的数字，辩护一方最有力的是刘易斯列出的机器工业（本书的概括）。巴西与墨西哥 1980 年以前的增长，与韩国从进口替代早早转向出口，都可以被两边引用（\chref{ch:eastasia}）。至今没有定论。
% 据: wd_41, wd_superior, wd_intermediate, wd_unclear, wd_korea, lw_protection

\subsubsection{大推进与援助}

第四个争论是协调与援助。按 Murphy、Shleifer 与 Vishny 的模型，协调失败可以让经济停在没有工业化的均衡里\cite{msv1988}。按他们自己的说法，多重均衡并不是自动出现的，即使金钱外部性重要，均衡也可以是唯一的\cite{msv1988}。苏联的经验在\chref{ch:soviet}，那里引的研究没有找到支持大推进理论的证据\cite{cggt2013}。按 Easterly 的检验，资金缺口模型几乎在每个国家都对不上，援助常常没有变成投资\cite{easterly1997}。主张推的一方说协调需要政府，怀疑的一方说政府的计划本身就是问题，印度的许可证与加纳的销售局都是例子（本书的概括）。至今没有定论。
% 据: ms_pecuniary, ms_unique, ea_tunisia, ea_60

\begin{deeper}[史料说明]
\textbf{版本。}刘易斯（1954）据互联网档案馆《曼彻斯特学派》第 22 卷的扫描本。普雷维什（1950）据互联网档案馆上的联合国文件 E/CN.12/89/Rev.1 扫描本，表 1 与第 12 页的数字按页面图像读。刘易斯与舒尔茨的演讲、自述与获奖页据诺贝尔奖网站。Murphy、Shleifer 与 Vishny 用的是 NBER 第 2708 号工作论文（1988 年 9 月），Easterly 用的是世界银行第 1807 号政策研究工作论文（1997 年 8 月），发表在期刊上的版本没有读到。印度的第二个计划草案大纲（1956）、第三个五年计划（1961）与加纳的七年计划大纲（1963）据互联网档案馆的缩微胶片扫描本。

\textbf{没有取得的。}辛格与赫希曼的著作、罗森斯坦—罗丹（1943）、罗斯托（1960）、纳克斯（1953）、舒尔茨（1964）、Krueger（1974）与 Little、Scitovsky 和 Scott（1970）都没有读到，罗森斯坦—罗丹与罗斯托只按别人的转述引。联合国 1949 年的价格文件没有取得。印度计划委员会网站上的第二个五年计划全文要求登录。

\textbf{文字层的误读。}普雷维什第 12 页的文字层把 36.5\% 读作 86.5\%，第 8 页把 1860 年代读作“I860's”。印度第三个计划的文字层把 11.5 与 8.5 读作“11:5”“8°5”，按页面图像读。

\textbf{不一致之处。}普雷维什第 8 页说 1930 年代只能买到 1860 年代的 63\%，表 1 却从 1876—1880 年算起，1936—1938 年的 64.1 对应的是下降 35.9\%。第 12 页说从 1870 年代到 1930 年代恶化 36.5\%，与 63\% 也不完全一致（本书的计算）。Easterly 引刘易斯说经济发展的中心事实是快速的资本积累，注的是文章首页第 139 页，这句话在第 155 页。国会图书馆罗斯托的记录写他 2003 年 2 月 14 日去世，GND 写 13 日。罗森斯坦—罗丹的生卒年只见于 GND。
\end{deeper}

\section{今天还剩什么}\label{sec:development-legacy}

\textbf{转折点。}劳动从农村转到城市，工资什么时候开始上升，仍是研究中国与其他快速工业化经济的一个问题，人口的变化在\chref{ch:demography}。刘易斯说，无限的劳动供给最终来自人口压力，是人口周期的一个阶段\cite{lewisbio1979}。
% 据: lb_demographic

\textbf{制度的长期影响。}资本只是起飞的一个条件。米塔制结束将近两百年后，按 Dell 的估计，受影响地区的家庭消费仍低约 25\%，儿童发育迟缓的比例高约 6 个百分点（\chref{ch:africa}）\cite{dell2010}。按 Dell 的解释，影响是通过土地占有与公共品的提供延续下来的\cite{dell2010}。技术与制度怎样进入增长理论，在\chref{ch:endogenous}。
% 据: dell_25, dell_6, dell_channels

\textbf{收入分配。}在刘易斯的模型里，增长的头一个阶段是利润的份额上升。增长与不平等的关系在\chref{ch:heterogeneity}。加纳的可可说明，销售局与汇率决定了农民拿到世界价格的多少。

\textbf{贸易与产业政策。}普雷维什的问题以新的形式回来。初级产品出口国的收入随价格起落，今天的关税与产业政策在\chref{ch:fragmentation}，新兴市场的危机在\chref{ch:emcrises}。

\textbf{援助与计划。}按 Easterly 的说法，哈罗德—多马模型早已过时，在援助的实践里却还活着\cite{easterly1997}。这个模型活下来，靠的是它与 40 年的数据相符吗？他的回答是否定的\cite{easterly1997}。成熟的计划经济与中国的大跃进及其饥荒在\chref{ch:planned}。
% 据: ea_ghost, ea_answer

\begin{takeaways}
\begin{itemize}
\item 刘易斯（1954）把经济分成资本主义部门与维持生计的部门。工资定在维持生计的收入之上约 30\%，利润再投资，资本主义部门扩大，利润与储蓄在国民收入中的份额上升，直到剩余劳动用完。
\item 普雷维什（1950）用英国的进出口价格说明，初级产品与工业品的价格之比从 1876—1880 年到 1936—1938 年降了约三分之一，主张工业化。有效保护率衡量关税对增加值的保护，可以远高于名义关税。
\item 大推进的模型里有两个均衡，单独建厂亏本，一起建厂都赚钱。资金缺口模型按 ICOR 算需要的投资，Easterly 的检验几乎在每个国家都否定了它。
\item 印度的计划者选择了重工业与更高的投资率，代价由消费承担；加纳的销售局压低可可的收购价，政府把差价用于别的用途。
\item 剩余劳动、贸易条件、进口替代与大推进四个问题，至今没有定论。
\end{itemize}
\end{takeaways}

\begin{selfcheck}
\begin{enumerate}
\item 按\exref{ex:development-lewis}的教学输入，资本主义部门雇 16 人时，全社会的储蓄率是多少？
\item 一种产品的名义关税 30\%，投入品按世界价格占产品价格的一半，投入品免税，有效保护率是多少？投入品征 30\% 呢？
\item 按式~\eqref{eq:development-push}，$L=100$，$a=1.5$，$v=0.2$，$F=15$ 时，没有工业化的均衡还存在吗？
\item 某国的 ICOR 为 3，目标增长率 6\%，国内储蓄率 12\%，资金缺口是多少？
\item 普雷维什表中的价格之比下降，出口国能买到的进口一定减少吗？
\item 为什么刘易斯说零边际产品对他的分析并不根本？
\end{enumerate}
\end{selfcheck}

\begin{answers}
\item $Y_c=31.2$，利润$10.4$，$Y=31.2+84=115.2$，储蓄率约 9.0\%。
\item $(0.3-0)/0.5=60\%$；$(0.3-0.5\times0.3)/0.5=30\%$，与名义关税相同。
\item 不存在。$20<1.2\times15=18$不成立，一家单独建厂也能赚钱，经济只有工业化的均衡。
\item $3\times6\%=18\%$，缺口$18\%-12\%=6\%$。
\item 不一定。收入贸易条件$(P_x/P_m)\,Q_x$还取决于出口量，出口量增长够快，能买到的进口仍会增加。
\item 模型要的是按一个固定的工资招得到人，劳动可以来自零工、家务与人口增长，农业劳动的边际产品是零还是很小，不改变这一点。
\end{answers}

回到 1952 年 8 月曼谷的那条路。刘易斯想到的答案，给了穷国一个提高储蓄率的机制，也说明这个机制先让利润增长，工资不涨。1957 年起，他做过加纳总理的经济顾问。到 1979 年，他说许多国家曾经向内，主要靠进口替代，后来几乎每个国家都发现了出口的“好处”\cite{lewis1979}。下一章讲的是最早转向出口的那几个经济。日本和东亚“四小龙”的高速增长，是生产率的奇迹，还是要素积累的结果？
% 据: ll_virtues
